% Options for packages loaded elsewhere
% Options for packages loaded elsewhere
\PassOptionsToPackage{unicode}{hyperref}
\PassOptionsToPackage{hyphens}{url}
%
\documentclass[
  11pt,
  letterpaper,
  DIV=11,
  numbers=noendperiod,
  twoside]{scrartcl}
\usepackage{xcolor}
\usepackage[left=1.1in, right=1in, top=0.8in, bottom=0.8in,
paperheight=9.5in, paperwidth=7in, includemp=TRUE, marginparwidth=0in,
marginparsep=0in]{geometry}
\usepackage{amsmath,amssymb}
\setcounter{secnumdepth}{3}
\usepackage{iftex}
\ifPDFTeX
  \usepackage[T1]{fontenc}
  \usepackage[utf8]{inputenc}
  \usepackage{textcomp} % provide euro and other symbols
\else % if luatex or xetex
  \usepackage{unicode-math} % this also loads fontspec
  \defaultfontfeatures{Scale=MatchLowercase}
  \defaultfontfeatures[\rmfamily]{Ligatures=TeX,Scale=1}
\fi
\usepackage{lmodern}
\ifPDFTeX\else
  % xetex/luatex font selection
  \setmathfont[]{Garamond-Math}
\fi
% Use upquote if available, for straight quotes in verbatim environments
\IfFileExists{upquote.sty}{\usepackage{upquote}}{}
\IfFileExists{microtype.sty}{% use microtype if available
  \usepackage[]{microtype}
  \UseMicrotypeSet[protrusion]{basicmath} % disable protrusion for tt fonts
}{}
\usepackage{setspace}
% Make \paragraph and \subparagraph free-standing
\makeatletter
\ifx\paragraph\undefined\else
  \let\oldparagraph\paragraph
  \renewcommand{\paragraph}{
    \@ifstar
      \xxxParagraphStar
      \xxxParagraphNoStar
  }
  \newcommand{\xxxParagraphStar}[1]{\oldparagraph*{#1}\mbox{}}
  \newcommand{\xxxParagraphNoStar}[1]{\oldparagraph{#1}\mbox{}}
\fi
\ifx\subparagraph\undefined\else
  \let\oldsubparagraph\subparagraph
  \renewcommand{\subparagraph}{
    \@ifstar
      \xxxSubParagraphStar
      \xxxSubParagraphNoStar
  }
  \newcommand{\xxxSubParagraphStar}[1]{\oldsubparagraph*{#1}\mbox{}}
  \newcommand{\xxxSubParagraphNoStar}[1]{\oldsubparagraph{#1}\mbox{}}
\fi
\makeatother


\usepackage{longtable,booktabs,array}
\newcounter{none} % for unnumbered tables
\usepackage{calc} % for calculating minipage widths
% Correct order of tables after \paragraph or \subparagraph
\usepackage{etoolbox}
\makeatletter
\patchcmd\longtable{\par}{\if@noskipsec\mbox{}\fi\par}{}{}
\makeatother
% Allow footnotes in longtable head/foot
\IfFileExists{footnotehyper.sty}{\usepackage{footnotehyper}}{\usepackage{footnote}}
\makesavenoteenv{longtable}
\usepackage{graphicx}
\makeatletter
\newsavebox\pandoc@box
\newcommand*\pandocbounded[1]{% scales image to fit in text height/width
  \sbox\pandoc@box{#1}%
  \Gscale@div\@tempa{\textheight}{\dimexpr\ht\pandoc@box+\dp\pandoc@box\relax}%
  \Gscale@div\@tempb{\linewidth}{\wd\pandoc@box}%
  \ifdim\@tempb\p@<\@tempa\p@\let\@tempa\@tempb\fi% select the smaller of both
  \ifdim\@tempa\p@<\p@\scalebox{\@tempa}{\usebox\pandoc@box}%
  \else\usebox{\pandoc@box}%
  \fi%
}
% Set default figure placement to htbp
\def\fps@figure{htbp}
\makeatother


% definitions for citeproc citations
\NewDocumentCommand\citeproctext{}{}
\NewDocumentCommand\citeproc{mm}{%
  \begingroup\def\citeproctext{#2}\cite{#1}\endgroup}
\makeatletter
 % allow citations to break across lines
 \let\@cite@ofmt\@firstofone
 % avoid brackets around text for \cite:
 \def\@biblabel#1{}
 \def\@cite#1#2{{#1\if@tempswa , #2\fi}}
\makeatother
\newlength{\cslhangindent}
\setlength{\cslhangindent}{1.5em}
\newlength{\csllabelwidth}
\setlength{\csllabelwidth}{3em}
\newenvironment{CSLReferences}[2] % #1 hanging-indent, #2 entry-spacing
 {\begin{list}{}{%
  \setlength{\itemindent}{0pt}
  \setlength{\leftmargin}{0pt}
  \setlength{\parsep}{0pt}
  % turn on hanging indent if param 1 is 1
  \ifodd #1
   \setlength{\leftmargin}{\cslhangindent}
   \setlength{\itemindent}{-1\cslhangindent}
  \fi
  % set entry spacing
  \setlength{\itemsep}{#2\baselineskip}}}
 {\end{list}}
\usepackage{calc}
\newcommand{\CSLBlock}[1]{\hfill\break\parbox[t]{\linewidth}{\strut\ignorespaces#1\strut}}
\newcommand{\CSLLeftMargin}[1]{\parbox[t]{\csllabelwidth}{\strut#1\strut}}
\newcommand{\CSLRightInline}[1]{\parbox[t]{\linewidth - \csllabelwidth}{\strut#1\strut}}
\newcommand{\CSLIndent}[1]{\hspace{\cslhangindent}#1}



\setlength{\emergencystretch}{3em} % prevent overfull lines

\providecommand{\tightlist}{%
  \setlength{\itemsep}{0pt}\setlength{\parskip}{0pt}}



 


\setlength\heavyrulewidth{0ex}
\setlength\lightrulewidth{0ex}
\usepackage[automark]{scrlayer-scrpage}
\clearpairofpagestyles
\cehead{
  Brian Weatherson
  }
\cohead{
  Reply to commentators
  }
\ohead{\bfseries \pagemark}
\cfoot{}
\makeatletter
\newcommand*\NoIndentAfterEnv[1]{%
  \AfterEndEnvironment{#1}{\par\@afterindentfalse\@afterheading}}
\makeatother
\NoIndentAfterEnv{itemize}
\NoIndentAfterEnv{enumerate}
\NoIndentAfterEnv{description}
\NoIndentAfterEnv{quote}
\NoIndentAfterEnv{equation}
\NoIndentAfterEnv{longtable}
\NoIndentAfterEnv{abstract}
\renewenvironment{abstract}
 {\vspace{-1.25cm}
 \quotation\small\noindent\emph{Abstract}:}
 {\endquotation}
\newfontfamily\tfont{EB Garamond}
\addtokomafont{disposition}{\rmfamily}
\addtokomafont{descriptionlabel}{\rmfamily}
\addtokomafont{title}{\normalfont\itshape}
\let\footnoterule\relax

\makeatletter
\renewcommand{\@maketitle}{%
  \newpage
  \null
  \vskip 2em%
  \begin{center}%
  \let \footnote \thanks
    {\itshape\huge\@title \par}%
    \vskip 0.5em%  % Reduced from default
    {\large
      \lineskip 0.3em%  % Reduced from default 0.5em
      \begin{tabular}[t]{c}%
        \@author
      \end{tabular}\par}%
    \vskip 0.5em%  % Reduced from default
    {\large \@date}%
  \end{center}%
  \par
  }
\makeatother
\RequirePackage{lettrine}

\renewenvironment{abstract}
 {\quotation\small\noindent\emph{Abstract}:}
 {\endquotation\vspace{-0.02cm}}

\setmainfont{EB Garamond Math}[
  BoldFont = {EB Garamond SemiBold},
  ItalicFont = {EB Garamond Italic},
  RawFeature = {+smcp},
]

\newfontfamily\scfont{EB Garamond Regular}[RawFeature=+smcp]
\renewcommand{\textsc}[1]{{\scfont #1}}

\renewcommand{\LettrineTextFont}{\scfont}
\usepackage{amsfonts}
\usepackage{lettrine}
\KOMAoption{captions}{tableheading}
\makeatletter
\@ifpackageloaded{caption}{}{\usepackage{caption}}
\AtBeginDocument{%
\ifdefined\contentsname
  \renewcommand*\contentsname{Table of contents}
\else
  \newcommand\contentsname{Table of contents}
\fi
\ifdefined\listfigurename
  \renewcommand*\listfigurename{List of Figures}
\else
  \newcommand\listfigurename{List of Figures}
\fi
\ifdefined\listtablename
  \renewcommand*\listtablename{List of Tables}
\else
  \newcommand\listtablename{List of Tables}
\fi
\ifdefined\figurename
  \renewcommand*\figurename{Figure}
\else
  \newcommand\figurename{Figure}
\fi
\ifdefined\tablename
  \renewcommand*\tablename{Table}
\else
  \newcommand\tablename{Table}
\fi
}
\@ifpackageloaded{float}{}{\usepackage{float}}
\floatstyle{ruled}
\@ifundefined{c@chapter}{\newfloat{codelisting}{h}{lop}}{\newfloat{codelisting}{h}{lop}[chapter]}
\floatname{codelisting}{Listing}
\newcommand*\listoflistings{\listof{codelisting}{List of Listings}}
\makeatother
\makeatletter
\makeatother
\makeatletter
\@ifpackageloaded{caption}{}{\usepackage{caption}}
\@ifpackageloaded{subcaption}{}{\usepackage{subcaption}}
\makeatother
\usepackage{bookmark}
\IfFileExists{xurl.sty}{\usepackage{xurl}}{} % add URL line breaks if available
\urlstyle{same}
% fallback for those not using the hyperref driver hyperxmp:
\makeatletter
\@ifundefined{xmpquote}{\newcommand{\xmpquote}[1]{#1}}{}
\makeatother
\hypersetup{
  pdftitle={Reply to commentators},
  pdfauthor={Brian Weatherson},
  pdfkeywords={\xmpquote{pragmatic
encroachment}, \xmpquote{interest-relativity}, \xmpquote{inquiry}, \xmpquote{checking}, \xmpquote{knowledge}, \xmpquote{sensitivity}},
  hidelinks,
  pdfcreator={LaTeX via pandoc}}


\title{Reply to commentators}
\author{Brian Weatherson}
\date{2026}
\begin{document}
\maketitle
\begin{abstract}
Matthew McGrath and Jane Friedman provided great comments on
\emph{Knowledge: A Human Interest Story}. Here I reply to them, making
some concessions, and relying on a distinction I wish I'd drawn in the
book. The distinction is between inquiries that aim at finding something
out, and inquiries that aim at checking something. These two inquiries
relate differently to knowledge. Any knowledge one has is fair game to
use as a starting point to finding things out, but there are different
constraints on how you check things. Making this distinction allows for
a cleaner statement of the role of knowledge in certain kinds of
inquiry, i.e., findings out, while allowing that knowers can sensibly
check what they already know.
\end{abstract}


\setstretch{1.1}
I'm really happy to reply to these comments, not least because
\emph{Knowledge: A Human Interest Story} (hereafter, KAHIS), wouldn't
exist without their work. Like many people, I started thinking about
interest-relativity because of Jeremy Fantl and Matthew McGrath
(\citeproc{ref-FantlMcGrath2002}{2002}), and the view defended in the
book is built on foundations that were laid in that paper and in their
(\citeproc{ref-FantlMcGrath2009}{2009}) book. And before reading what
Jane Friedman (\citeproc{ref-Friedman2019a}{2019b},
\citeproc{ref-Friedman2019b}{2019a}, \citeproc{ref-Friedman2020}{2020})
had to say about inquiry, I wouldn't have thought to restructure the
theory in an inquiry-centric way.

I'll start with McGrath's comments, taking them slightly out of order,
because what I say in response to one of his later questions helps set
up what I want to say both to his earlier questions, and to Friedman's
comments.

\section{Reply to McGrath}\label{sec-mcgrath}

\subsection{Checking and Inquiring}\label{sec-mcgrath-checking}

I'll start with McGrath's third comment. In KAHIS I agreed that
typically one cannot inquire into what one knows, but argued that there
are some atypical cases where you can. In particular, when the purpose
of inquiry is getting sensitive beliefs, it is fine to inquire into
things one knows. Since saying to oneself ``\emph{p}, therefore
\emph{p}'' would not increase the sensitivity of one's belief in
\emph{p}, this isn't a good inquiry, even though it starts with what one
knows and uses clearly valid steps. (Is this a counterexample to the
principle that you can always start with what you know? No, because
what's wrong with this inquiry is not where it starts, but the result
that running through it fails at the aim of increasing sensitivity.)

McGrath suggests that these cases are better described as
\emph{checking} than as inquiring. As he points out, ordinary language
agrees here. When he looks into the office at the insistence of a
sceptical colleague, it is natural to say he is \emph{checking} whether
the (department) chair is there, and unnatural to say one is inquiring
into whether she is there.

I'm happy to accept his amendment, but I don't think this is the end of
the story, for a couple of reasons. First, this leaves open what exactly
it is to check something. It isn't just inquiring into what we know.
Grabbing an old flashlight to take to the basement while a storm
approaches, my spouse might sensibly say ``You should check that's still
working.'' This isn't because we already know it is working; quite the
contrary. So checking is compatible with ignorance. On the other hand,
not all inquiry is checking. If I go into a philosopher's archives and
start reading with no fixed goal, I might be trying to find out
something about them, but I'm not checking anything. Exactly what
checking is feels like a bigger question.\footnote{The most important
  recent work on checking is \emph{Knowing and Checking}
  (\citeproc{ref-Melchior2019}{Melchior 2019}). We'll come back to it at
  more length in Section~\ref{sec-friedman-checking}.}

Or, perhaps I should say, it feels like \emph{two} bigger questions. One
question is how the ordinary term `check' gets used. The other question
is whether there is a theoretically interesting class of activities
picked out by this word, or in the vicinity of what this word picks out.
I suspect the answer to the latter question is \emph{yes}, but that's
for another day.

This leads to my other worry about McGrath's argument. I said above that
it isn't natural to describe McGrath's checkers as inquiring into
something. But it's rarely natural to say someone is `inquiring into'
something. That's a theorist's construction, not ordinary language.
There's a response to this, and it's right there in what McGrath said.
Some inquirers are \emph{trying to find something out}. The person who
checks because of a sceptical colleague is not doing that. They know
where the chair is, they aren't trying to find it out.

This all suggests a nice synthesis. There are two related activities:
\emph{checking} and \emph{finding out}. We can do both at once; I find
out whether the flashlight is working by checking it. But they can come
apart in both directions. Open-ended archival research is finding out
without checking; satisfying a sceptical colleague is checking without
finding out. Which of these is `inquiry'? I'd rather say they
\emph{both} are kinds of inquiry; more on this in
Section~\ref{sec-friedman-checking}.

What's more important is that with this distinction in hand, we can say
something more systematic about the relationship between knowledge and
inquiry. When you are simply trying to find things out, it's permissible
to start with any (relevant) fact you know. That's enough to get the
argument for pragmatic encroachment off the ground. More carefully, I
want to accept the principle KS-FIND, and reject the principle
KS-CHECK.\footnote{The `KS' here is short for `Knowledge Starts'.}

\begin{description}
\tightlist
\item[KS-FIND]
When you are trying to find something out, it's permissible to start
with anything relevant you know.
\item[KS-CHECK]
When you are trying to check something, it's permissible to start with
anything relevant you know.
\end{description}

If you have that pair of judgments, it makes sense to accept KP-FIND,
and reject KP-CHECK.\footnote{The `KP' here is short for `Knowledge
  Permits'.}

\begin{description}
\tightlist
\item[KP-FIND]
It is not permissible, by the standards of inquiry, to try to find out
something you already know.
\item[KP-CHECK]
It is not permissible, by the standards of inquiry, to try to check
something you already know.
\end{description}

My guess is that McGrath and I agree about all of these verdicts, but we
might disagree about which activities are \emph{inquiries}. Let's leave
that last point for now, because it will make more sense to discuss it
in the context of Friedman's comments.

\subsection{Questions About Probability}\label{sec-mcgrath-probability}

The view in KAHIS has a striking consequence concerning probability. It
says that when one is trying to find out how likely it is that \emph{p},
either the answer is 1, or one doesn't know \emph{p}. McGrath says that
this is implausible when \emph{p} is \emph{At least one player will hit
a free throw in tonight's NBA games}. (Assume it's mid-season, and there
are ten games on.) If I stick to this verdict, I have to say,
implausibly according to McGrath, that either the answer to the
likelihood question is 1, or the inquirer doesn't know \emph{p}. I'm not
going to try wriggling out of this; I think the second option is right:
the inquirer doesn't know \emph{p}. That is not exactly what intuition
says, so I better argue for it.

Start with how you'd try to figure out this likelihood. You'd look at
the frequency of made free throws over the last few hundred games, the
number of free throw attempts awarded per game, perhaps something about
how those attempts are distributed. What you would not do, under any
circumstances, is start with ``Well, at least one will be made tonight,
so let's take that as given.'' That premise is simply unavailable in
\emph{this} inquiry, because of which question is being asked.

In the previous subsection I distinguished checking whether \emph{p}
from trying to find out whether \emph{p}, and I endorsed a strong
principle about the latter: when you are trying to find something out,
you can use whatever relevant fact you know. I'd like to keep that
principle exceptionless, and the view that you lose knowledge that
\emph{p} when working out the likelihood of \emph{p} is the only way to
keep it. On McGrath's picture, the inquirer knows that some player will
make a free throw, that's obviously relevant to the likelihood they are
trying to deduce, but they can't use it. The connection between
knowledge and sensible methods for finding stuff out has been broken,
and once it is broken, for me the interest of inquiring into the nature
of knowledge goes with it.

McGrath objects that the reasoning I'm attributing to the knower is
invalid: the premise \emph{p} is about the world, the conclusion
\emph{Pr(p) = 1} is about an information state. True, but compare:
\emph{This player hits 90\% of his free throws, so the probability that
he'll hit this one is 0.9}. That also goes from world to information
state, it's also invalid (as can be seen by the fact that it fails if
you add extra information), yet it looks perfectly fine. The natural
thing to say is that it's enthymematic; the suppressed premise is
something like \emph{and that's all the relevant information I have}.
The reasoning from \emph{p} is enthymematic in the same way, with the
suppressed premise even weaker: just that \emph{p} is among what can be
taken as given. There are tricky questions here about exactly what we're
doing when we reason from worldly facts to probabilities, but I suspect
however they get answered, there will be a way to endorse the reasoning
I'm attributing to the knower.

McGrath suggests the naturalness of the following speech raises a
problem for my view here. One could concede a minuscule probability that
no free throw is made, then add ``But that's not going to happen;
someone will make one''? That's fine I think, but look what it does. It
\emph{sets a possibility aside}. Contra Lewis
(\citeproc{ref-Lewis1996}{1996}), I don't think the fact that we are
attending to a possibility means we have to take it seriously. We can
propose that we're not going to consider some possibility,\footnote{Here
  I'm basically in agreement with the criticisms that Blome-Tillmann
  (\citeproc{ref-MBT2009a}{2009}) makes of the Rule of Attention that
  Lewis (\citeproc{ref-Lewis1996}{1996}) proposes} and provided that
everyone else goes along with this, we can act as if we know it. If this
ignoring makes sense given our interests, we may in fact know it, so
it's not just acting as if. All I insist on is that in doing that we are
tacitly agreeing that we're no longer working out how many 9's there are
at the start of the probability that at least one free throw will be hit
tonight.

\subsection{Causal and Constitutive
Dependence}\label{sec-mcgrath-constitutive}

I can be brief about the fourth point. McGrath grants that the Das-style
argument I make in chapter 7 is correct as far as it goes: every
plausible theory makes knowledge counterfactually depend on
truth-irrelevant factors. That's already enough to push back against
some of the claims that pragmatic encroachment is a radically new view.
But he suggests that there is a way in which it is new, and which might
be behind the critics' intuition that it is an unwanted intruder. The
other dependencies are \emph{causal}, while the dependence of knowledge
on what one can properly take for granted is \emph{constitutive}. He
suggests, rather generously to our mutual critics, that what's radical
about pragmatic encroachment is that it makes knowledge constitutively
depend on non-truth-relevant factors. Here's where I get off the boat; I
think that constitutive dependence is already there once we insist on
environmental reliability.

Consider the analogy McGrath uses. He agrees that reliability
considerations mean that knowledge might depend on interests because
they make me more likely to go to some places rather than others. But he
insists this making is causal, ``like weighting a coin.'' Let's think
about that coin. Weighting the coin causes something: a mass
distribution. The relationship between the mass distribution and the
chance, however, is not causal; it's constitutive. There's no causal
pathway from the mass distribution changing to the chance changing, and
no temporal gap for that pathway to play out in. The same goes for the
mental case. Suppose I prefer ebooks to printed books, and so am more
likely to read something in an ebook. Whatever caused the preference
caused a mental state, and that state, together with background
conditions, grounds the likelihood facts. The preference doesn't sit
causally upstream of the likelihood. There's no chain running between
them. Like with the coin, there's no time for that chain to run through.
As soon as I have the preference, I'm more likely to read ebooks. The
preference might cause my actually reading the ebook not the printed
book, but it doesn't cause my being more likely to do so.

So a non-truth-relevant fact, that I prefer ebooks to printed books,
constitutively fixes the environment that's relevant to reliability
evaluation. Facts about that environment fix, again constitutively, how
likely I am to come across certain errors. So if we hold fixed the
external world, which contains veridical printed books and erroneous
ebooks, my change in preference is linked constitutively, not merely
causally, to a decrease in reliability. Put another way, Das's real
point was about individuation. What counts as \emph{my} environment,
which reference class \emph{my} reliability is assessed against, depends
on my interests. This individuation is not a causal matter at all. So
constitutive dependence of knowledge on truth-irrelevant factors is not
what separates pragmatic encroachment from traditional theories, as long
as versions of reliabilism count as traditional. It may still be that
encroachment strikes people as radical; but if so, the explanation is
not that it alone has something truth-irrelevant among the things
knowledge constitutively depends on, because on inspection, many many
people do.

\subsection{Odds and Stakes}\label{sec-mcgrath-odds}

I'll end with McGrath's second point, where I actually agree with most
of what he says. In KAHIS, I say that high stakes aren't directly
relevant to knowledge changes. McGrath is right, though, to say that for
some propositions it will be very hard in practice to get knowledge to
change without high stakes.

McGrath's example of this is \emph{The Battle of Hastings was in 1066}.
Funnily enough, when I was drafting KAHIS, I had originally used that
example for the blue sentence. I dropped it because I wasn't sure it
really was irrational to play Blue-True in that case! The intuition
about rational play is load-bearing in the argument, and I didn't think
in this case it would bear the load.

This points to a broader methodological point. I'm certainly not wanting
to sign on for a general principle about logical strength in the
red-blue game. If the blue sentence is \emph{2+2 is less than 100} and
the red is \emph{2+2=4}, I don't think either play is irrational. Maybe
I could be convinced that one is, there are some really deep questions
about the status of weak dominance arguments lurking in the background
here, but for now I'm happy to sign up for either being rational
here.\footnote{One might worry here that the appeal to weak dominance I
  made in the précis is inconsistent with what I say here. The careful
  version of what I'm saying is that if \(X\) weakly dominates \(Y\),
  that entails that choosing \(X\) in a two-way choice between \(X\) and
  \(Y\) is \emph{permissible}, but does not entail that it is
  \emph{required}.} A nice bonus of this is that I avoid the sceptical
worry McGrath presses, that maybe sometimes I'll end up saying one
doesn't know that 2+2=4.

Still, even if I agree that it would take really high stakes to lose
knowledge in the Hastings case, I want to insist that this isn't enough
to resuscitate a connection between stakes and knowledge. If one loses
knowledge in these cases, it's because the odds change. It might be
practically impossible for the odds to change without the stakes
changing, but that just means high stakes situations are tightly
correlated with knowledge loss, not that being in a high stakes
situation directly explains the knowledge loss.

The reason for this is that I'm worried that the notion of stakes isn't
sharp enough to play a central role in grounding knowledge. Anderson and
Hawthorne (\citeproc{ref-AndersonHawthorne2019a}{2019}) argue that once
we look at cases with more than two options, it's really unclear what
the stakes are, and no precisification of the vague informal notion can
work. Rather than trying to answer the challenge they posed, my view
simply dodges it. `High stakes' has a role in informally characterising
the theory, not in the most careful statement of it.

\section{Reply to Friedman}\label{sec-friedman}

\subsection{How Interests Matter}\label{sec-friedman-interests}

Jane Friedman notices a couple of tensions within KAHIS, and asks for
some clarifications on each. These are excellent points to press on.
I'll clarify one thing, and concede one thing. Let's start with the
clarification.

When I'm discussing practical choices, like what move to make in games,
I stress the importance of the principle (\(\star\)), where \(L\)
measures how much one might lose if \emph{p} is false, and \(G\)
measures how much one might gain if \emph{p} is true.

\[\frac{\Pr(p)}{1-\Pr(p)} > \frac{L}{G} \tag{$\star$}\]

That's definitive of when interests matter for a very particular kind of
inquiry: an inquiry into what to do when there are only two (salient)
options, and the salient rule is \emph{Maximise expected utility!}
Friedman notes that this principle, which starts looking so central,
doesn't do any work when I start talking about theoretical interests.
That's a fair point. Indeed, an even more general point would also be
fair; (\(\star\)) doesn't do much work in three-way choices, or in
choices like the ones discussed in chapter 6 where considerations other
than utility are relevant.

The short answer is, as Friedman suggests, the core principle in the
book is that knowledge is a constraint on starting points in inquiry.
There is one special kind of inquiry, inquiry into which of two options
to choose when only utility is relevant to which is appropriate, where
the main constraint on whether \emph{p} is a proper starting point is
(\(\star\)). But even in that case, it's not that there is some basic
principle of knowledge that implies (\(\star\)) is important; it's
derived from the general principle about proper starting points.

So why do interests matter? They matter because the right principle
connecting knowledge and inquiry quantifies over inquiries one is
\emph{or should be} conducting. And sometimes, one's interests determine
what one should be inquiring into. Those interests could be either
practical or theoretical interests: presumably if someone is in fact
inquiring into something they have to be at least minimally interested
in it. Conversely, if one is really interested in something, there will
be situations where one should be inquiring into it.

One of these theoretical investigations does a bit of work in KAHIS: the
case where Anisa merely imagines playing the red-blue game. Friedman
says she likes this case but isn't sure what intuitions she has there,
and is a little doubtful that it tells us a lot. Actually, I don't have
any strong intuitions there either; if anything my intuition is the
opposite of the view in KAHIS that she loses historical knowledge even
once she starts this hypothetical inquiry. But I was never in the
business of trying to maximise agreement with case-intuitions. The
following two things strike me as the most plausible things to say about
it. First, there is no intuition about the case strong enough to
overturn the plausibility of applying KS-FIND here. Second, it is a bit
funny for her to make the following speech: ``Hypothetically, if the
blue sentence was \emph{Agincourt was in 1415}, I could take either
option. Oh look, that is the blue sentence. Guess my only option is
red.'' If she doesn't lose knowledge in the hypothetical, theoretical,
inquiry, I'm not sure what blocks that speech.

This case shows that even hypothetical options matter for knowledge.
Both practical interests and theoretical interests are relevant. In
KAHIS I suggested that the term `interest' in the last sentence might
usefully unify things here. The thought was that it wasn't a coincidence
that we use a similar word to describe something practical, like (1),
and something theoretical, like (2).

\begin{enumerate}
\def\labelenumi{(\arabic{enumi})}
\tightlist
\item
  How vaccines are distributed around here is relevant to my (practical)
  interests.
\item
  I'm (theoretically) interested in how vaccines are distributed around
  here.
\end{enumerate}

On reflection, I'm worried this is a bit of a coincidence. I'm not sure
you do get the same word form used for `interest' in these claims in
enough languages. You do, I think, see similar words used when you
translate (1) and (2) into most \emph{European} languages, but maybe not
beyond that. So I'm a bit worried that calling my view
`interest-relative' because it takes the kind of facts reported by both
(1) and (2) to be relevant to knowledge, because they are relevant to
proper starting points for inquiry, is almost a pun. Friedman suggests
that when I'm discussing `theoretical interests' it might be better to
describe the view as inquiry-relative than interest-relative. It might
do better to unify things to say the whole view is inquiry-relative,
it's just that sometimes the inquiry is into what to do.

\subsection{Finding Out and Checking}\label{sec-friedman-checking}

This is the point where I think I need to make a concession, and it
relates to what I said in Section~\ref{sec-mcgrath-checking}. The
problem is that I'd talked myself into a muddle. Here are two things I
say in KAHIS.

In the red-blue game, the player loses historical knowledge. They must
lose that knowledge, because if they still had it, they could reason
their way into rationally playing Blue-True, and that's impossible.

In the double-checking case, the player doesn't lose knowledge. They
have knowledge, but given their aim of increasing sensitivity, they
can't use it. That implies there are further constraints, given their
aims, on what they can start inquiry with beyond knowledge and
relevance.

This suggests a question, one that Friedman presses. Why can't the
invariantist say the same thing about game playing as I say about
checking? The player retains their historical knowledge, but they can't
use it because using it is incompatible with their aim of maximising
utility. What stops there being some other aim? Good question!

In KAHIS I gave some arguments that there couldn't be such a reason;
that was the point of the arguments from Moore-paradoxicality. But those
seemed like arguments from symptoms; they didn't explain why the cases
are different. I now have the sketch of an answer; namely the principles
I set out at the end of Section~\ref{sec-mcgrath-checking}.

The difference between the cases is just the difference between finding
out and checking. The game player is trying to find out what to do. Or,
if that over-intellectualises it, they are trying to do whichever thing
is the answer to the question \emph{What to do?} The double-checker is,
well, checking. The relationship between knowledge and these two
activities is rather different: KS-FIND and KP-FIND are true; KS-CHECK
and KP-CHECK are false.

Combined with some intuitive judgments about who is checking and who is
finding out, this gets plausible verdicts in all the cases. Brown's
surgeon and Florian the coffee maker are checking, and so don't lose
knowledge; Darja who is trying to decide whether to check, the inquirer
who bootstraps, and the red-blue game player are finding out (or at
least trying to), and do lose knowledge.

In the previous paragraph, I said Brown's surgeon is a checking case.
One might worry this is question-begging. As Friedman points out, there
are two distinct choices here. The first is \emph{whether} to check or
go straight into operating, the second is \emph{how} to check. Even if
you grant me that it's consistent to say that she knows which kidney to
remove but can't use that knowledge in checking, there is still the
first decision. Here it looks like (\(\star\)) tells us that she loses
knowledge. Indeed, in a similar case in KAHIS, the case of Darja, I said
it was a knowledge loss. Why is the surgeon different?

The difference, I think, is that the surgeon is a rule-follower, not an
expected utility maximiser. She checks because she has, sensibly,
adopted the policy of always checking. If she doesn't check, she
violates a rule, and that's ruled out. That rule violation is not
conditional on which kidney is meant to be removed. So I think
ultimately the surgeon isn't actually making the first inquiry; it's
already settled by her adoption of a standing policy. It doesn't matter
here whether she adopted that policy or it is required by her job,
though the point that not checking has costs independent of which kidney
is to be removed is even cleaner if it is a required step. If, like
Darja, the surgeon is making a real decision about whether to check or
go straight to operating, and the only cost of operating is removing the
wrong kidney, then I am committed to saying that she either loses
knowledge or is being irrational. That doesn't strike me as the most
realistic version of the case; the point of these pre-surgery checklists
is to have policy replace decision-making in stressful circumstances.

At this point I'm leaning pretty heavily on distinctive features of
checking. What does this imply for the relationship between knowledge
and \emph{inquiry}? I think to answer that it helps to separate two
questions:

\begin{itemize}
\tightlist
\item
  Is checking a form of inquiry?
\item
  Does checking require losing knowledge?
\end{itemize}

McGrath argues the answer to the first is `no'. I'm inclined to say it
is `yes', for a few reasons. One is a blatant appeal to authority: it's
what Jane Friedman says in `Checking Again': ``First, checking is
inquiring.'' (\citeproc{ref-Friedman2019b}{Friedman 2019a, 85}). A
second is that some checks, like checking the flashlight, are
uncontroversially inquiries. A third is that checks, even into what one
knows, involve the same movements as these uncontroversial inquiries. A
fourth (or maybe a defence of the third), is that the best example of
someone who acts like an inquirer but is actually not---the murderous
Morse in ``Inquiry and Belief'' (\citeproc{ref-Friedman2019a}{Friedman
2019b})---is also not a checker.\footnote{My intuition is that neither
  the case where Morse actually killed the doctor, nor the case where he
  falsely believes that he did, are cases of checking, but what's really
  crucial is the case where he actually did it.}

I won't rehearse the arguments that the answer to the second is `no',
because that's what I did in chapter 5 of KAHIS, and they aren't that
different to what Arianna Falbo (\citeproc{ref-Falbo2021}{2021}) and
Elise Woodard (\citeproc{ref-Woodard2021}{2024}) have already offered.
So I can locate my position in Table~\ref{tbl-checking}, where the rows
are the first question, and the columns are the second question.

\begin{longtable}[]{@{}
  >{\raggedright\arraybackslash}p{(\linewidth - 4\tabcolsep) * \real{0.3768}}
  >{\raggedright\arraybackslash}p{(\linewidth - 4\tabcolsep) * \real{0.3188}}
  >{\raggedright\arraybackslash}p{(\linewidth - 4\tabcolsep) * \real{0.3043}}@{}}
\caption{Four options about checking, knowledge, and
inquiry}\label{tbl-checking}\tabularnewline
\toprule\noalign{}
\begin{minipage}[b]{\linewidth}\raggedright
\end{minipage} & \begin{minipage}[b]{\linewidth}\raggedright
Loses knowledge: YES
\end{minipage} & \begin{minipage}[b]{\linewidth}\raggedright
Loses knowledge: NO
\end{minipage} \\
\midrule\noalign{}
\endfirsthead
\toprule\noalign{}
\begin{minipage}[b]{\linewidth}\raggedright
\end{minipage} & \begin{minipage}[b]{\linewidth}\raggedright
Loses knowledge: YES
\end{minipage} & \begin{minipage}[b]{\linewidth}\raggedright
Loses knowledge: NO
\end{minipage} \\
\midrule\noalign{}
\endhead
\bottomrule\noalign{}
\endlastfoot
\textbf{Checking is inquiry: YES} & Friedman & Melchior \\
\textbf{Checking is inquiry: NO} & ??? & McGrath \\
\end{longtable}

I'm in the top-right corner, along with Guido Melchior.\footnote{See
  Melchior (\citeproc{ref-Melchior2019}{2019}), and more recently
  Melchior (\citeproc{ref-Melchior2023}{2023}), especially the replies
  to Miščević and to Siscoe.} Friedman's two 2019 papers commit her to
the top-left corner, while McGrath's comments here put him in the
bottom-right corner. I don't know who, if anyone, is in the bottom-left
corner. It seems undermotivated, but now I've identified a position in
logical space I'm optimistic it will get filled.

Relying on this distinction between finding out and checking is, I
think, a concession. In KAHIS I suggested that the core principles about
knowledge I was using followed from general principles about knowledge
and inquiry, principles that had to be carefully qualified to not crash
against the checking cases. But I didn't explain why there wouldn't be
similar qualifications in the game-playing cases, or in cases like
working out the probability that at least one free throw will be scored.
What I should have said was that when someone is trying to find
something out, whether that's when something happened, how likely
something is, or what to do, they can start with any relevant knowledge.
But when they are trying to check something, another important kind of
inquiry, knowledge doesn't have a privileged role.

\section*{References}\label{references}
\addcontentsline{toc}{section}{References}

\protect\phantomsection\label{refs}
\begin{CSLReferences}{1}{0}
\bibitem[\citeproctext]{ref-AndersonHawthorne2019a}
Anderson, Charity, and John Hawthorne. 2019. {``Knowledge, Practical
Adequacy, and Stakes.''} \emph{Oxford Studies in Epistemology} 6:
234--57.

\bibitem[\citeproctext]{ref-MBT2009a}
Blome-Tillmann, Michael. 2009. {``Knowledge and Presuppositions.''}
\emph{Mind} 118 (470): 241--94. doi:
\href{https://doi.org/10.1093/mind/fzp032}{10.1093/mind/fzp032}.

\bibitem[\citeproctext]{ref-Falbo2021}
Falbo, Arianna. 2021. {``Inquiry and Confirmation.''} \emph{Analysis} 81
(4): 622--31. doi:
\href{https://doi.org/10.1093/analys/anab037}{10.1093/analys/anab037}.

\bibitem[\citeproctext]{ref-FantlMcGrath2002}
Fantl, Jeremy, and Matthew McGrath. 2002. {``Evidence, Pragmatics, and
Justification.''} \emph{Philosophical Review} 111 (1): 67--94. doi:
\href{https://doi.org/10.2307/3182570}{10.2307/3182570}.

\bibitem[\citeproctext]{ref-FantlMcGrath2009}
---------. 2009. \emph{Knowledge in an Uncertain World}. Oxford: Oxford
University Press.

\bibitem[\citeproctext]{ref-Friedman2019b}
Friedman, Jane. 2019a. {``Checking Again.''} \emph{Philosophical Issues}
29 (1): 84--96. doi:
\href{https://doi.org/10.1111/phis.12141}{10.1111/phis.12141}.

\bibitem[\citeproctext]{ref-Friedman2019a}
---------. 2019b. {``Inquiry and Belief.''} \emph{No{û}s} 53 (2):
296--315. doi:
\href{https://doi.org/10.1111/nous.12222}{10.1111/nous.12222}.

\bibitem[\citeproctext]{ref-Friedman2020}
---------. 2020. {``The Epistemic and the Zetetic.''}
\emph{Philosophical Review} 129 (4): 501--36. doi:
\href{https://doi.org/10.1215/00318108-8540918}{10.1215/00318108-8540918}.

\bibitem[\citeproctext]{ref-Lewis1996}
Lewis, David. 1996. {``Desire as Belief {II}.''} \emph{Mind} 105 (418):
303--13. doi:
\href{https://doi.org/10.1093/mind/105.418.303}{10.1093/mind/105.418.303}.

\bibitem[\citeproctext]{ref-Melchior2019}
Melchior, Guido. 2019. \emph{Knowing and Checking: An Epistemological
Investigation}. Routledge Studies in Epistemology. New York: Routledge.
doi:
\href{https://doi.org/10.4324/9780429030239}{10.4324/9780429030239}.

\bibitem[\citeproctext]{ref-Melchior2023}
---------. 2023. {``Replies to the Critics of \emph{Knowing and
Checking: An Epistemological Investigation}.''} \emph{Acta Analytica} 38
(1): 95--131. doi:
\href{https://doi.org/10.1007/s12136-022-00541-0}{10.1007/s12136-022-00541-0}.

\bibitem[\citeproctext]{ref-Woodard2021}
Woodard, Elise. 2024. {``Why Double-Check.''} \emph{Episteme} 21:
664--67. doi:
\href{https://doi.org/10.1017/epi.2022.22}{10.1017/epi.2022.22}.

\end{CSLReferences}



\noindent Draft.


\end{document}
